%% Hello World Tree — TeX adapter.
%%
%% TeX is the emit language. Every other adapter in this repository reads an artifact; this one writes
%% one. That asymmetry is deliberate and it is what closes the loop:
%%
%%   decode   artifact -> structure   (the other seventeen)
%%   encode   structure -> artifact   (this one, plus SLEIGH's mirrored engine)
%%
%% A vessel that can only read is half a vessel. The moment you can emit the structure back as
%% something a typesetting engine understands, you can produce a document — and a document is the one
%% artifact form that survives being printed, archived, and read by a person who has no tooling at all.
%%
%% The macros here generate that document. They take the same node/edge shape every adapter emits, so a
%% skeleton from any language can be rendered to a PDF with no intermediate format.

\documentclass[10pt]{article}

\usepackage[T1]{fontenc}
\usepackage[utf8]{inputenc}
\usepackage{xcolor}
\usepackage{tikz}
\usepackage{geometry}
\usepackage{hyperref}

\geometry{letterpaper, margin=0.9in}

%% ── The palette ─────────────────────────────────────────────────────────────
%% The same values the web and the QML shell use, so a printed vessel is recognisably the same object.
%% Dioxazine is far too dark for print, so it is approximated with a desaturated plum; the hum and the
%% gold survive unchanged because they are the identity.

\definecolor{void}{HTML}{030208}
\definecolor{diox}{HTML}{1A0A28}
\definecolor{hum}{HTML}{E5007E}
\definecolor{humhot}{HTML}{FF2D9B}
\definecolor{indigo}{HTML}{2E3BFF}
\definecolor{gold}{HTML}{FFD400}
\definecolor{typehi}{HTML}{FFFFFF}
\definecolor{plum}{HTML}{7A2350}

%% ── Node and edge tables ────────────────────────────────────────────────────
%% The adapter's input. Populated by the emitter before the document is rendered, so a single invocation
%% produces a complete artifact from a complete graph.
%%
%% Nodes carry no weight. A dead vessel has no data, so there is nothing to weight, and a weight field
%% would invite values to be invented for it.

\newcount\hwtnodecount
\hwtnodecount=0

\newcount\hwtedgecount
\hwtedgecount=0

%% \hwtaddnode{id}{kind}{label}
\def\hwtaddnode#1#2#3{%
  \global\advance\hwtnodecount by 1
  \expandafter\gdef\csname hwtnode@\the\hwtnodecount @id\endcsname{#1}%
  \expandafter\gdef\csname hwtnode@\the\hwtnodecount @kind\endcsname{#2}%
  \expandafter\gdef\csname hwtnode@\the\hwtnodecount @label\endcsname{#3}%
}

%% \hwtaddedge{from}{to}{kind}
\def\hwtaddedge#1#2#3{%
  \global\advance\hwtedgecount by 1
  \expandafter\gdef\csname hwtedge@\the\hwtedgecount @from\endcsname{#1}%
  \expandafter\gdef\csname hwtedge@\the\hwtedgecount @to\endcsname{#2}%
  \expandafter\gdef\csname hwtedge@\the\hwtedgecount @kind\endcsname{#3}%
}

%% ── Node styling ────────────────────────────────────────────────────────────
%% Kind determines the mark. Size and colour, in that order of importance: a reader distinguishes
%% differently-sized marks before they distinguish colours, which is why colour is the secondary channel
%% and not the primary one.

\def\hwtkindcolor#1{%
  \def\tmp{#1}%
  \ifx\tmp\hwtnodefile gold%
  \else\ifx\tmp\hwtnodeimport indigo%
  \else\ifx\tmp\hwtnodeinstruction plum%
  \else humhot%
  \fi\fi\fi
}

\def\hwtnodefile{file}
\def\hwtnodeimport{import}
\def\hwtnodeinstruction{instruction}

%% \hwdrawnode{x}{y}{kind}{label}
\def\hwdrawnode#1#2#3#4{%
  \begin{scope}[shift={(#1,#2)}]
    %% halo
    \fill[\hwtkindcolor{#3}, opacity=0.18] (0,0) circle (0.62);
    %% core
    \fill[\hwtkindcolor{#3}] (0,0) circle (0.22);
    %% label
    \node[anchor=west, font=\ttfamily\tiny, text=black!80] at (0.34,0) {#4};
  \end{scope}
}

%% ── The document ────────────────────────────────────────────────────────────

\begin{document}

\pagecolor{void}
\color{typehi}

\begin{center}
  {\ttfamily\Large\bfseries HELLO WORLD TREE}\\[4pt]
  {\ttfamily\small\color{gold} A VESSEL, NOT YET FILLED}
\end{center}

\vspace{6pt}

%% The honest statement. On a dead vessel this is the single most important line in the document,
%% because everything below it is structure without data and a reader needs to know that before they
%% read anything into the shape.
\begin{center}
  {\ttfamily\footnotesize\color{humhot}%
   \the\hwtnodecount\ nodes \quad \the\hwtedgecount\ edges \quad no flow}
\end{center}

\vspace{10pt}
\hrule
\vspace{10pt}

%% ── The graph ───────────────────────────────────────────────────────────────
%% Radial layout computed here rather than passed in, because the emitter should not have to know about
%% paper size. Each node gets an angle from its index; nodes are placed on a ring sized to the count.
%%
%% The result is not a good layout for a large graph. It is a correct layout for a small one, and at
%% this stage every vessel is small.

\begin{center}
\begin{tikzpicture}[scale=0.86, transform shape]

  \def\ringradius{4.6}
  \pgfmathsetmacro{\step}{360/\the\hwtnodecount}

  %% edges first, so nodes draw over them
  \foreach \i in {1,...,\the\hwtedgecount} {%
    \pgfmathsetmacro{\fx}{0}%
    \pgfmathsetmacro{\fy}{0}%
    \pgfmathsetmacro{\tx}{0}%
    \pgfmathsetmacro{\ty}{0}%
    %% A real emitter resolves node ids to coordinates before this point. The loop structure is here
    %% so the shape of the document is visible; the resolution is the emitter's job, not TeX's.
    \draw[indigo, opacity=0.22, line width=0.4pt] (\fx,\fy) -- (\tx,\ty);
  }

  %% nodes
  \foreach \i in {1,...,\the\hwtnodecount} {%
    \pgfmathsetmacro{\angle}{\i*\step}
    \pgfmathsetmacro{\nx}{\ringradius*cos(\angle)}
    \pgfmathsetmacro{\ny}{\ringradius*sin(\angle)}
    \hwdrawnode{\nx}{\ny}%
      {\csname hwtnode@\i @kind\endcsname}%
      {\csname hwtnode@\i @label\endcsname}
  }

\end{tikzpicture}
\end{center}

\vspace{10pt}
\hrule
\vspace{10pt}

%% ── The structure, as text ──────────────────────────────────────────────────
%% A printed graph is lossy: overlap hides nodes, and colour is unavailable in a photocopy. So the same
%% structure is also emitted as a table, which is the form that survives being printed, scanned, and
%% read by something that is not a person.
%%
%% This redundancy is deliberate. Every artifact the vessel emits should be legible in at least two
%% channels, because any single channel fails eventually.

\begin{center}
{\ttfamily\footnotesize
\begin{tabular}{@{}lll@{}}
  \textbf{\color{gold}kind} & \textbf{\color{gold}label} & \textbf{\color{gold}id} \\
  \hline
  \foreach \i in {1,...,\the\hwtnodecount} {%
    \csname hwtnode@\i @kind\endcsname &
    \csname hwtnode@\i @label\endcsname &
    \csname hwtnode@\i @id\endcsname \\
  }
\end{tabular}
}
\end{center}

\vfill

\begin{center}
  {\ttfamily\tiny\color{plum}%
   artifact in, proof out \quad · \quad reproducible evidence outranks documentation}
\end{center}

\end{document}
